\begin{statement}{B.24}{The unified integral class-two local-order law}{Proved}
For the cubic at $\ell=3$ and the plus quartic at
$\ell\in\{3,5\}$, the image of the actual class-two section
set in its abelian input space is the subset described below.
In every case the rule has the same form: every displayed
local correction is integral, and its minimum coordinate
valuation tends to $+\infty$ outside finite sets of primes.
Valuation of zero means $+\infty$.

\emph{Cubic input and its forty-four scalar corrections.}
Use $O=C_0$, $K=\mathbb Q(\zeta_3)$, and the ordered basis
\[
 (f_1,f_2,A_0,A_1,A_2,B_0,B_1,B_2,C_1,C_2)
                                                               \tag{B.24.1}
\]
of the actual abelian lattice $V=T_3J_{E,D}$. Here the last
eight symbols denote positive peripheral cocharacters, and
\[
 A_i=[0:-\zeta_3^i:1],\quad
 B_i=[-\zeta_3^i:0:1],\quad C_i=[1:-\zeta_3^i:0].
\]
Write $W$ for their span and let $U_1,U_2\in W$ be the
columns of
\[
 U=\begin{pmatrix}
 0&-1\\2&-2\\1&0\\0&-2\\1&-1\\2&0\\1&-2\\2&-1
 \end{pmatrix},\qquad
 R_3=\begin{pmatrix}0&-1\\1&-1\end{pmatrix}.
                                                               \tag{B.24.2}
\]
The compact inclusion has block column $\binom{3I_2}{U}$.
Put $w=(3f_1+U_1)\wedge(3f_2+U_2)$, whose coefficient at
$A_0\wedge A_2$ is one.

Let $\alpha\in H^1(\mathbb Q,V)$. For a prime
$\mathfrak q\nmid3$ of $K$, put $q=N\mathfrak q$.
Its elliptic Frobenius is
$A_E=aI_2+bR_3$, where $a+b\zeta_3$ is the primary
Jacobi-sum generator, congruent to one modulo three.
At an inert rational prime $p$ it is $-p$ and $q=p^2$.
The actual affine matrix is
\[
 A_{\mathfrak q}=
 \begin{pmatrix}
 A_E&0\\ U(A_E-qI_2)/3&qI_8
 \end{pmatrix}.                                                \tag{B.24.3}
\]
Take a tame generator $t_{\mathfrak q}$ and arithmetic
Frobenius $F_{\mathfrak q}$, with $FtF^{-1}=t^q$.
Set $a_{\mathfrak q}=\alpha(F_{\mathfrak q})$ and
$v_{\mathfrak q}=\alpha(t_{\mathfrak q})
 =\sum_d b_{\mathfrak q,d}e_d\in W$.
Let $k_{\mathfrak q,d}$ be an integral representative of
the unramified Kummer class of the following actual framed
point of $J_{E,D}$: $z_d=\mathcal O(d-O)$ with frames
$s_d^{-1}$ at $d$, $s_O$ at $O$, and one elsewhere.
The parameters are $s_{A_i}=x/z$, $s_{B_i}=y/z$ and
$s_{C_i}=z/x$. All the finite-field input is given by
\[
 \begin{split}
 3z_d&=\left(
   \operatorname{lead}_r(L_d/L_O)/
   \operatorname{lead}_O(L_d/L_O)\right)_{r\ne O},\\
 L_{A_i}&=y+\zeta_3^iz,\quad
 L_{B_i}=x+\zeta_3^iz,\quad
 L_{C_i}=x+\zeta_3^{2i}y,\quad L_O=x+y.
 \end{split}                                                    \tag{B.24.4}
\]
At $d$ the leading value divides by $s_d^3$, and at
$O$ it multiplies by $s_O^3$; elsewhere it is evaluation.
The Kummer class is read from divisions of $\overline z_d$
in $J_{E,D}(\overline{\mathbb F}_q)$ by applying $F-1$.
It lies in the finite quotient $V/(A_{\mathfrak q}-1)V$.

Form, in the basis (B.24.1),
\[
 \begin{split}
 D_{\mathfrak q}
 &=a_{\mathfrak q}\wedge v_{\mathfrak q}
       -\sum_d b_{\mathfrak q,d} k_{\mathfrak q,d}\wedge e_d\\
 &=f_1\wedge X_1+f_2\wedge X_2+Y,
       \qquad X_i\in W,\quad Y\in\Lambda^2W .
 \end{split}                                                    \tag{B.24.5}
\]
Regard $X=(X_1,X_2)^{\rm t}$ as a two-by-eight matrix.
The following operations take place over $\mathbb Q_3$:
\[
 \begin{split}
 Z&=(A_E-I_2)^{-1}X,\qquad Z=(Z_1,Z_2)^{\rm t},\\
 Y_*&=\frac{Y-\tfrac13(U_1\wedge X_1+U_2\wedge X_2)}{q-1}
          +\tfrac13(U_1\wedge Z_1+U_2\wedge Z_2),\\
 z_*&=(Y_*)_{A_0\wedge A_2}.
 \end{split}                                                    \tag{B.24.6}
\]
Define $\rho_{\mathfrak q}(\alpha)$ to be the minimum
$3$-adic valuation of these forty-four coordinates:
\[
 \begin{array}{c|c}
 \text{coordinate block}&\text{correction}\\ \hline
 f_1\wedge f_2&9z_*\\
 f_1\wedge W&-Z_1+3z_*U_2\\
 f_2\wedge W&-Z_2-3z_*U_1\\
 \Lambda^2W\text{ except }A_0\wedge A_2&
                   -Y_*+z_*U_1\wedge U_2 .
 \end{array}                                                    \tag{B.24.7}
\]
The complete cubic image is
\[
 \left\{\alpha:\ 
   \rho_{\mathfrak q}(\alpha)\geq0\text{ for all }\mathfrak q\nmid3,
   \quad \rho_{\mathfrak q}(\alpha)\longrightarrow+\infty\right\}.
                                                               \tag{B.24.8}
\]
This is a two-by-two formula for the full intrinsic
class-two quotient, with no omission of the integral
index-nine compact gluing.

\emph{Quartic input and its scalar order functions.}
Use the plus quartic $x^4+y^4+z^4=0$ and
$F=\mathbb Q(\zeta_8)$. Let $\alpha_i=\zeta_8^{1+2i}$,
and take cusps
$A_i=[0:\alpha_i:1]$, $B_i=[\alpha_i:0:1]$,
$C_i=[\alpha_i:1:0]$, with $O=B_0$.
Use unit parameters $x/z,y/z,z/y$, respectively.
The abelian section origin is the average of the four
conjugate $B_i$ tangents. It has an actual class-two lift.
In the splitting of B.23, every input has unique coordinates
\[
 (b,\mathbf r)\in H^1(\mathbb Q,W)\oplus R^3,\qquad
 e_{\mathbf r}=\sum_{i=1}^3r_i j_i h,
                                                               \tag{B.24.9}
\]
where $h$ is any fixed generator of $H^1(\mathbb Q,V)=R$
and $j_i:V\to A$ are the actual compact inclusions.
Equivalently the first factor is
\[
 H^1(\mathbb Q,W)=
 (\widehat{F^\times}_{\,\ell})^2
 \times\ker\left(N_{F/\mathbb Q}:
      \widehat{F^\times}_{\,\ell}\to
      \widehat{\mathbb Q^\times}_{\,\ell}\right).
                                                               \tag{B.24.10}
\]

Restrict $b$ to $F$ in the order
$(A_0,A_1,A_2,A_3,B_1,B_2,B_3,C_0,C_1,C_2,C_3)$.
Let $\mathbf 1$ denote the length-four all-one column and put
\[
 \begin{gathered}
 H=\begin{pmatrix}
 -8&-3&-2&-3\\-3&-8&-3&-2\\
 -2&-3&-8&-3\\-3&-2&-3&-8
 \end{pmatrix},\quad r=\begin{pmatrix}-5\\-6\\-5\end{pmatrix},\\
 B=\begin{pmatrix}-10&-6&-4\\-6&-12&-6\\-4&-6&-10\end{pmatrix},
 \qquad
 N=\begin{pmatrix}
 H&{\bf1}r^{\rm t}&-4{\bf1}{\bf1}^{\rm t}\\
 r{\bf1}^{\rm t}&B&r{\bf1}^{\rm t}\\
 -4{\bf1}{\bf1}^{\rm t}&{\bf1}r^{\rm t}&H
 \end{pmatrix},\\
 v=(-\tfrac12,-\tfrac12,-\tfrac12,-\tfrac12,
       -\tfrac58,-\tfrac34,-\tfrac58,
       -\tfrac12,-\tfrac12,-\tfrac12,-\tfrac12)^{\rm t}.
 \end{gathered}                                                 \tag{B.24.11}
\]
For every $i<j$ let $k=\kappa_\ell(2)$ and define
\[
 p_{ij}=v_i-N_{ij}/8,\quad q_{ij}=v_j-N_{ij}/8,\quad
 X_{ij}=b_i+p_{ij}k,\quad Y_{ij}=b_j+q_{ij}k .
                                                               \tag{B.24.12}
\]
This gives fifty-five pairs of explicit Kummer classes.

Here is their scalar test on all completed inputs.
Take $\epsilon=1+\sqrt2$ and a generator $\pi_{\mathfrak p}$
of every prime ideal of $F$. Every Kummer class has unique
coordinates
\[
 x=a_\epsilon\kappa_\ell(\epsilon)
       +\sum_{\mathfrak p}a_{\mathfrak p}
                         \kappa_\ell(\pi_{\mathfrak p}),
 \qquad a_{\mathfrak p}\longrightarrow0.
                                                               \tag{B.24.13}
\]
For a pair $X,Y$ denote its two coefficient families by
$a,b$. At $\mathfrak q\nmid\ell$, put
$Q=N\mathfrak q$, $e=\nu_\ell(Q-1)$, and choose
a primitive $\ell^e$-th root $\eta_{\mathfrak q}$
in its residue field. For $u=\epsilon$ or
$u=\pi_{\mathfrak p}$ with $\mathfrak p\ne\mathfrak q$,
let $r_{\mathfrak q,u}$ be the integer in $[0,\ell^e)$
determined by
\[
 \overline u^{(Q-1)/\ell^e}
       =\eta_{\mathfrak q}^{\,r_{\mathfrak q,u}}.
\]
Set
\[
 \begin{split}
 A_{\mathfrak q}
 &=a_\epsilon r_{\mathfrak q,\epsilon}
    +\sum_{\mathfrak p\ne\mathfrak q}
                    a_{\mathfrak p}r_{\mathfrak q,\pi_{\mathfrak p}},\\
 B_{\mathfrak q}
 &=b_\epsilon r_{\mathfrak q,\epsilon}
    +\sum_{\mathfrak p\ne\mathfrak q}
                    b_{\mathfrak p}r_{\mathfrak q,\pi_{\mathfrak p}},\\
 \rho^{ij}_{\mathfrak q}(b)
 &=\nu_\ell(A_{\mathfrak q}b_{\mathfrak q}
                -B_{\mathfrak q}a_{\mathfrak q})-e .
 \end{split}                                                    \tag{B.24.14}
\]
If $e=0$ take all $r_{\mathfrak q,u}=0$ and
$\rho^{ij}_{\mathfrak q}=+\infty$.
The toric test is
\[
 \rho^{ij}_{\mathfrak q}(b)\geq0\text{ for every }\mathfrak q\nmid\ell,
 \qquad
 \rho^{ij}_{\mathfrak q}(b)\longrightarrow+\infty
          \quad\text{for all }i<j.
                                                               \tag{B.24.15}
\]

There is one further order function for the mixed part.
At a rational $p\notin S$, put $v_p=b(t_p)\in W$ and
\[
 \overline u_p=
  (F_{E,p}-I)^{-1}h(F_p)\pmod V,\qquad
 \operatorname{ord}(\overline u_p)=\ell^{m_p}.
                                                               \tag{B.24.16}
\]
The class is in $(V\otimes\mathbb Q_\ell/V)^{F_p}$,
identified with $E(\mathbb F_p)[\ell^\infty]$. In particular
$m_p\leq\nu_\ell(\#E(\mathbb F_p))$. For $\mathbf r\ne0$
write $t=\min_i\nu_\ell(r_i)$ and set
\[
 \rho^{\rm mix}_p(b,\mathbf r)
  =\nu_\ell(v_p)+t-m_p\quad(m_p>0),\qquad
 \rho^{\rm mix}_p=+\infty\quad(m_p=0).
                                                               \tag{B.24.17}
\]
For $\mathbf r=0$ set every $\rho^{\rm mix}_p=+\infty$.
The mixed test is
\[
 \rho^{\rm mix}_p(b,\mathbf r)\geq0\text{ for all }p\notin S,
 \qquad
 \rho^{\rm mix}_p(b,\mathbf r)\longrightarrow+\infty .
                                                               \tag{B.24.18}
\]
The complete quartic image consists exactly of the pairs
$(b,\mathbf r)$ satisfying both (B.24.15) and (B.24.18).
Thus the function cutting out the image is fully specified
by scalar residue logarithms and the orders of the actual
local Kummer classes (B.24.16). Replacing $m_p$ by the
possibly larger valuation of the group order would change
the criterion.

All tests are independent of cochain representatives,
compatible path choices, and Frobenius prolongations.
The finite logarithms in (B.24.14) are used modulo their
specified powers; they are not asserted to equal full
$\ell$-adic Frobenius values. Every nonempty fibre is a
torsor under the actual central group
$H^1_{\rm cts}(\mathbb Q,\Lambda^2 A)$, with $A=V$
in the cubic case.
\end{statement}
\begin{proof}
\emph{The continuous localization mechanism.}
For a finite free unramified coefficient $T$ on a fixed
$S$-integer base, suppose
$H^2(G_S,T/\ell^n)=H^3(G_S,T/\ell^n)=0$ and the finite-base
$H^1$ is finite at every $n$. First take the direct limit
of finite-coefficient localization over finite sets of
removed primes. The support complex has two terms
\[
 \bigoplus_{\mathfrak q\notin S} T/\ell^n
   \ \xrightarrow{\ q^{-1}F_{T,\mathfrak q}-1\ }\
 \bigoplus_{\mathfrak q\notin S} T/\ell^n .
\]
Both systems of terms have surjective reduction maps.
Their derived inverse limit is therefore the literal
two-term complex on
$\mathcal R(T)=\{(z_{\mathfrak q}):z_{\mathfrak q}\in T,\
z_{\mathfrak q}\to0\}$. The finite-base vanishings and
finiteness kill the adjacent derived limits. Thus
\[
 H^2_{\rm cts}(G,T)=
 \mathcal R(T)/
 (q^{-1}F_{T,\mathfrak q}-1)_{\mathfrak q}\mathcal R(T).
                                                               \tag{B.24.19}
\]
The image is its literal image, not its closure.
These are the localization and inverse-system comparison
maps of \citet[II.0.1, II.1.7, II.2.9]{BMilne2006}
and \citet[Proposition 1.6, Theorem 2.2]{BJannsen1988}.
This proves the common origin of integrality and decay
in the two cases.

\emph{Cubic geometry and the inverse formula.}
The nine cusps are $E[3]$. With $A_0,A_2$ as the ordered
three-torsion basis their Weil pairing is $\zeta_3$.
Their coordinates in this basis are
$(1,0),(2,2),(0,1),(2,0),(1,1),(0,2),(2,1),(1,2)$.
Pulling back the generalized-Jacobian extension by $[3]$
trivializes its characters $\mathcal O(d-O)$.
The vector frame transforms by the inverse Weil pairing,
so its gluing row is $(v_d,-u_d)$, giving (B.24.2).
Equivariance of the columns $\binom{3I_2}{U}$ gives
(B.24.3). The normalized tangent-line divisors give
(B.24.4). At the two exceptional evaluations the leading
values are, respectively,
\[
 \begin{array}{c|cc}
 d&\operatorname{lead}_d(L_d/L_O)&
           \operatorname{lead}_O(L_d/L_O)\\ \hline
 A_i&1/3&3\\
 B_i&1/3&-3\\
 C_i,\ i=1,2&-1/(3(1-\zeta_3^i))&
                       -3(1-\zeta_3^{2i}).
 \end{array}
\]
They are units away from three. Passing from the endpoint
frame of a tangent path to this divisor frame introduces
only signs of order two, whose pro-three Kummer torsor is
canonically trivial. Thus $k_d$ is the actual abelian
path cocycle, with its stated sign.

The good smooth pair and its unit tangent give unramified
based pro-three action outside three. The relevant
prime-to-residue-characteristic specialization, including
inertia identification, is
\citet[Expos\'e XIII, \S\S2.9--2.12]{BSGA1}.
The tame relation then forces the inertia of $\alpha$
into $W$: its elliptic projection would be an eigenvector
of $F_E$ with eigenvalue $q$, which does not exist.
In class-two BCH coordinates the group law is
$(x,z)(y,t)=(x+y,z+t+x\wedge y/2)$.
On a peripheral generator the actual arithmetic linear
term is
$D_g(e_d)=-\chi(g)k_d(g)\wedge e_d$.
Compute $FtF^{-1}=t^q$ in the pullback extension over
$\alpha$. The central Frobenius coordinate cancels and
the equation for its inertia coordinate $m$ is
\[
 (q^{-1}\Lambda^2 A_{\mathfrak q}-1)m=-D_{\mathfrak q}.
                                                               \tag{B.24.20}
\]
This identifies the actual support representative, not
just its image in a finite quotient.

The kernel of the operator in (B.24.20) is the primitive
line $\mathbb Z_3w$. All its other eigenvalues are
nonzero by their Frobenius weights. Set
$g_i=f_i+U_i/3$. In this rational basis the operator is
$A_E-I_2$ on $g_i\wedge W$ and $q-1$ on $\Lambda^2W$.
Writing (B.24.5) in that basis and inverting gives
$m_0=\sum_i f_i\wedge Z_i+Y_*$ for the positive
right-hand side $D_{\mathfrak q}$.
Subtracting $z_*w$ eliminates the pivot $A_0\wedge A_2$.
Negating this representative gives exactly (B.24.7).
Thus these are the integral coordinates of the unique
solution in $\Lambda^2V/\mathbb Z_3w$, not those of an
unjustified direct sum with the compact line.
The determinant of the resulting forty-four-dimensional
map, up to a unit, is
$\#E(\mathbb F_q)^8(q-1)^{28}/9$.
The factor nine is retained by the $9z_*$ term.
The script independently verifies both inverse identities
on all forty-four basis columns for six distinct elliptic
Frobenius matrices, in exact rational arithmetic.

For clarity the finite-base cubic vanishing has a direct
Kummer proof. The ring $\mathbb Z[\zeta_3]$ is Euclidean,
and there is exactly one prime above three.
Thus $\operatorname{Pic}(\mathcal O_K[1/3])=0$ and
its Brauer group is zero: a Brauer class can have
only one finite invariant, whose sum must be zero.
As $\mu_3\subset K$, Kummer theory gives
$H^2(G_{K,\{3\}},\mathbb F_3)=0$.
The matrices (B.24.3) are block unipotent modulo three,
so the residual action has three-group image.
A composition series therefore
gives the vanishings needed in (B.24.19) for
$T=\Lambda^2V$ and every exponent.
Continuity gives $v_{\mathfrak q}\to0$, and all other
entries of (B.24.5) are integral and bounded. Hence its
support family is restricted. The fixed integral
pivot section of $\Lambda^2V/\mathbb Z_3w$ lifts
restricted families to restricted families.
Equation (B.24.19) now proves both directions of
(B.24.8). Restriction from $\mathbb Q$ to $K$ is
injective on the obstruction group, since its
corestriction composite is multiplication by two.

\emph{Quartic tangent origin and all peripheral terms.}
The tangent functions are
$L_{A_i}=y-\alpha_i z$, $L_{B_i}=x-\alpha_i z$,
$L_{C_i}=x-\alpha_i y$.
The ratios $L_d/L_O$ have divisors $4d-4O$. At a cusp $d$, the leading
coefficient of $L_d$ in its specified parameter is
$\alpha_i/4$. Evaluating all ratios of leading terms
of $L_d/L_O$ gives
\[
 4z_d=(\zeta_8^{K_{rd}}(\sqrt2)^{N_{rd}})_{r\ne O}.
                                                               \tag{B.24.21}
\]
The independently recomputed matrix $N$ is exactly
(B.24.11); the script checks all 121 equalities in
$\mathbb Q[t]/(t^4+1)$. Roots of unity have zero
odd-primary Kummer class, so
$\kappa(z_d)=N_{\bullet d}\kappa_\ell(2)/8$.
This fixes every path contribution integrally.
If $L(e_d)=e_d\wedge N_{\bullet d}/8$, the actual
restricted action on $W\oplus\Lambda^2W$ is
\[
 g(x,z)=\bigl(\chi(g)x,\
        \chi(g)^2z+\chi(g)k(g)L(x)\bigr).
                                                               \tag{B.24.22}
\]
The average of the four conjugate $B_i$ tangent classes
has abelian coordinate $k v$; summing the three nonzero
columns in (B.24.21) and dividing by four gives
the stated $v$. Its central coordinate
$k(g)^2L(v)/2$ satisfies the class-two cocycle identity
in (B.24.22). This proves the existence of its
class-two lift over $F$.
Here is the exact averaging construction used to descend.
In an extension with abelian kernel $T$, choose four
coset representatives $g_i$ for $G/H$, lifts $e_i$,
and a splitting $s_H$ over $H$.
For an element $e$ above $g$, write
$gg_i=g_{gi}h_i(g)$ and set
\[
 B(e)=\sum_i
       e\,e_i\,s_H(h_i(g))^{-1}e_{gi}^{-1}\ \in T .
\]
Multiplying the factors and reindexing gives
$B(ef)=B(e)+gB(f)$ and $B(z)=4z$ for $z\in T$.
Hence $(-B(e)/4)e$ is an independent-of-lift section.
Apply this first to the abelian extension; its restricted
class is the displayed average $kv$.
Apply it a second time to the pulled-back central
extension with its just-constructed $H$-splitting.
This gives an actual class-two lift over $\mathbb Q$
preserving that abelian origin. No nonabelian averaging
has been used.
Both constructions are defined over $G_{\mathbb Q,S}$.
Indeed the pair has good reduction outside two and
the tangent leading constants are units there;
the same specialization result as above applies at
every $p\notin S$. Locally the image of arithmetic
inertia in the coefficientwise tangential splitting
centralizes the full free pro-$\ell$ geometric group $N$.
Different such splittings have the same inertia image:
their ratio lies in $Z(N)=1$.
The closed subgroup generated by these inertia images
is therefore normal in the full arithmetic extension,
centralizes $N$, and intersects $N$ trivially.
Its arithmetic image is exactly the kernel of
$G_{\mathbb Q}\to G_{\mathbb Q,S}$.
Quotienting by it gives the required extension over
$G_{\mathbb Q,S}$ with unchanged geometric kernel.

In the shifted origin, direct substitution into the BCH
cocycle identity gives for the $(i,j)$ peripheral
obstruction
$X_{ij}\cup Y_{ij}+d\Phi_{ij}$, where
\[
 \Phi_{ij}(g)=
 \tfrac12 b_i(g)b_j(g)+q_{ij}k(g)b_i(g)
                 +\tfrac12 p_{ij}q_{ij}k(g)^2.
                                                               \tag{B.24.23}
\]
This equality is checked as a polynomial identity in
the two cocycle values and $\chi(g)$ in the script.
The coefficient splitting and Shapiro's cocycle maps
give (B.24.10). Class number one and
$\mathcal O_F^\times=\mu_8\langle1+\sqrt2\rangle$
give the coordinates (B.24.13); the unit assertion
is also checked using the certified fundamental
units of the same quartic field.

For $R(2)$ the localization differential is $Q-1$.
The tame relation for $X\cup Y$ gives its support
representative
$A_{\mathfrak q}b_{\mathfrak q}
 -B_{\mathfrak q}a_{\mathfrak q}$, with an immaterial
common sign. The finite residue logarithms determine
each unramified Kummer value modulo $\ell^e$.
Replacing the full values by these representatives
changes its correction by an integral family tending
to zero: the differences are in
$(Q-1)\mathbb Z_\ell$, and they are multiplied by
$a_{\mathfrak q},b_{\mathfrak q}\to0$.
All infinite sums converge because their coefficients
tend to zero and the integer residues are integral.
B.23 gives the finite-base $R(2)$ vanishing. Applying
(B.24.19) therefore proves precisely (B.24.15), with
its decay condition. Restriction/corestriction of
degree four returns the answer to $\mathbb Q$.

\emph{Why no unevaluated compact action remains in the law.}
Write the actual linear class-two action as
$g(x,z)=(\rho_A(g)x,\rho_B(g)z+D_gx)$ and
$c_g=D_g\rho_A(g)^{-1}$, a cocycle in
$\operatorname{Hom}(A,\Lambda^2A)$.
For each Tate projection
$\lambda:\Lambda^2A\to L$, $L=R(1)$ or
$R(1)\chi_{-1}$, the perfect pairings
\[
 B_\lambda(x,y)=\langle x,y\rangle
       \quad\text{or}\quad \langle Ix,y\rangle
                                                               \tag{B.24.24}
\]
identify $\operatorname{Hom}(V,L)$ with $V$.
The first is alternating; the second is symmetric,
because $I$ is symplectic and $I^2=-1$.
Consequently $[\lambda c j_i]$ is a scalar multiple
of $[h]$, by B.23. Its evaluation on $g h(t)$ is a
scalar multiple of the corresponding self-cup, modulo
an explicit coboundary.

Both relevant self-cups vanish. The symmetric one
has primitive $-B_\lambda(h,h)/2$.
For the alternating one, at $2$ every point of
$E(\mathbb Q_2)$ reduces into the nonsingular
subgroup of order two: an integral odd $x$ would make
$x(x^2+1)$ have odd valuation one, and a nonintegral
point reduces to infinity. The formal group kernel
is pro-two, so $E(\mathbb Q_2)[\ell^n]=0$.
Local duality and the prime-to-two Euler characteristic
\citep[I, Theorem 2.1 and Lemma 2.9]{BMilne2006}
give $H^1(\mathbb Q_2,E[\ell^n])=0$.
At good primes other than $\ell$, $h$ is unramified,
so its finite cup factors through a procyclic group
of $\ell$-cohomological dimension one. The real
contribution is zero. The Brauer reciprocity sequence
then kills the only remaining contribution, at $\ell$
\citep[VIII, Theorem 4.2]{BMilneCFT2020}.
The full-field Kummer $H^1$ transitions are onto,
so the continuous alternating cup also vanishes.
The injection (B.23.4) gives its vanishing on $G_S$.

All other coefficient projections have zero finite-base
$H^2$ by B.23. Thus every pure compact input
$e_{\mathbf r}$ has an actual continuous class-two
lift over $G_S$, including its linear arithmetic
term $D_g e_{\mathbf r}(t)$.
This argument neither assigns a value to the Ceresa
class nor assumes it to be zero.
For a pure toric input, restriction to $F$ puts the
entire obstruction in $B_W$ by (B.24.22), so
its other projections are boundaries by degree-four
restriction/corestriction.
Finally the full obstruction for $b+e_{\mathbf r}$
differs from the sum of the two pure obstructions
and the ordered mixed cup
$b(g)\wedge g e_{\mathbf r}(t)$ by
$d(b\wedge e_{\mathbf r})/2$.
After subtracting the just-constructed primitives,
the only remaining terms lie in the two distinct
summands $B_W$ and $M$.

\emph{The mixed local order and its sufficiency.}
Apply (B.24.19) to $M$, using B.23.
The tame relation in the ordered-cup extension has
central defect
$-p\,v_p\wedge e_{\mathbf r}(F_p)$.
Dividing by $p$ in the support map, its correcting
equation is
\[
 (p^{-1}F_{M,p}-1)c_p
        =v_p\wedge e_{\mathbf r}(F_p).
                                                               \tag{B.24.25}
\]
Write $F_{W,p}=pP_p$. The inertia cocycle relation
gives $P_pv_p=v_p$, so the unique rational solution is
\[
 c_p=v_p\wedge
          j_{\mathbf r}(F_{E,p}-I)^{-1}h(F_p).
                                                               \tag{B.24.26}
\]
Uniqueness follows also from the Frobenius weights:
every eigenvalue of $p^{-1}F_M$ has absolute value
$\sqrt p$, so none is one.
The boundary map in
$0\to V\to V\otimes\mathbb Q_\ell
       \to V\otimes\mathbb Q_\ell/V\to0$
identifies (B.24.16) with the unramified local
Kummer class of $h$. Choose a fractional representative
of $\overline u_p$ with coordinates in
$\mathbb Z[1/\ell]\cap[0,1)$.
It has minimum valuation $-m_p$ if $m_p>0$,
and is zero otherwise. The change from the actual
inverse in (B.24.26) is integral; multiplied by
$v_p\to0$, it is an integral restricted family.
In the integral tensor lattice
$M=W\otimes R^3\otimes V$ the valuation of a nonzero
pure tensor is the sum of its factor valuations.
This gives exactly (B.24.17)--(B.24.18).

For explicit finite evaluation of $m_p$, it suffices
to compute the marked cocycle $h$ and elliptic
Frobenius modulo $\ell^a$ with
$a=\nu_\ell(\#E(\mathbb F_p))$.
No uncomputed infinite precision is required.
Compatible representatives $h_n$ can, if needed,
be constructed by enumerating finite Galois
extensions unramified outside $S$ containing
$\mathbb Q(E[\ell^n])$, solving their finite cocycle
equations, and choosing a primitive class lifting
$h_{n-1}$. Such an enumeration is effective by
defining polynomials and discriminant tests.
B.23 proves that the required primitive class
exists and that the coefficient maps are onto,
so each finite stage terminates.
This prescription retains the actual representation
and Frobenius marking; it does not replace $h$
by an arbitrary homomorphism.

If all tests hold, the correcting families belong
to $\mathcal R$ in (B.24.19), which provides an
actual continuous primitive. One may also construct
it exponent by exponent: the prescribed corrections
have finite support modulo each $\ell^n$;
localization supplies a primitive with those
inertia values. The remaining choices form a finite
torsor under unramified one-cocycles.
Every finite collection of compatibility conditions
can be met at its largest exponent, so their
inverse limit supplies a continuous primitive.
The boundary corrections above return an actual
class-two section. Conversely any section gives
these same corrections by its tame relation.
This proves necessity as well as sufficiency.
Changes of representatives alter the corrections
by integral restricted families, as the displayed
local equations show. Finally $H^0(\mathbb Q,A)=0$
by Frobenius weights. A difference of two aligned
central section coordinates is a one-cocycle,
and central conjugacy is its ordinary coboundary.
This proves the fibre assertion.
\end{proof}

